%% !TeX root = main.tex
% =========================
%         PAQUETES
% =========================
\usepackage[utf8]{inputenc}
\usepackage[T1]{fontenc}
\usepackage[spanish,es-nodecimaldot]{babel}
\usepackage{geometry}
\geometry{left=2.0cm,right=2.0cm,top=2.7cm,bottom=2.8cm,headheight=15pt}
\usepackage{amsmath,amssymb,amsfonts,amsthm,mathtools}
\usepackage{mathrsfs}
\usepackage{bm}
\usepackage{physics}
\usepackage{braket}
\usepackage{graphicx}
\usepackage{float}
\usepackage{xcolor}
% \usepackage{microtype} % desactivado para evitar advertencias de compatibilidad en Overleaf
\usepackage{tikz}
\usepackage{enumitem}
\usepackage{array}
\usepackage{booktabs}
\usepackage{hyperref}
\usepackage{subfiles}
\usepackage{fancyhdr}
\usepackage{titlesec}
\usepackage{setspace}
\usepackage{csquotes}
\usepackage{tensor}
\usepackage{lastpage}
\usepackage{tcolorbox}
\tcbuselibrary{breakable,skins}
\usepackage{tikz}
\usetikzlibrary{arrows.meta,calc,decorations.pathreplacing,positioning,shapes.geometric}
\usepackage{pgfplots}
\pgfplotsset{compat=1.18}


% Estilos TikZ usados en figuras didácticas del apunte
\tikzset{
    eje/.style={-{Latex[length=2.2mm,width=1.5mm]}, thick, gray!70!black},
    vecazul/.style={-{Latex[length=2.8mm,width=1.8mm]}, very thick, UdeCAzul},
    vecdorado/.style={-{Latex[length=2.8mm,width=1.8mm]}, very thick, UdeCDorado},
    vecverde/.style={-{Latex[length=2.8mm,width=1.8mm]}, very thick, green!50!black},
    frente/.style={fill=UdeCAzul!7, draw=UdeCAzul!75!black, thick},
    cajadiagrama/.style={rounded corners=2mm, draw=UdeCAzul!45, fill=UdeCAzul!5, inner sep=4pt, align=center}
}

% =========================
%       PALETA VISUAL
% =========================
\definecolor{UdeCAzul}{HTML}{223C6A}
\definecolor{UdeCDorado}{HTML}{E69B0A}
\definecolor{UdeCGris}{HTML}{A0A0A0}
\definecolor{UdeCGrisClaro}{HTML}{F4F5F7}
\graphicspath{{./}{../}{chapters/}}

% =========================
%     CONFIGURACIÓN GENERAL
% =========================
\setstretch{1.08}
\setlength{\parskip}{0.35em}
\setlength{\parindent}{0pt}

\hypersetup{
    colorlinks=true,
    linkcolor=UdeCAzul,
    citecolor=UdeCAzul,
    urlcolor=UdeCAzul,
    pdftitle={Electrodinámica 2},
    pdfauthor={José Rosas Sepúlveda}
}


\titleformat{\chapter}[display]
  {\Huge\bfseries\color{UdeCAzul}}
  {\chaptername\ \thechapter}
  {0.6em}
  {}

\titleformat{\section}{\Large\bfseries\color{UdeCAzul}}{\thesection}{0.6em}{}
\titleformat{\subsection}{\large\bfseries\color{UdeCAzul}}{\thesubsection}{0.6em}{}
\titleformat{\subsubsection}{\normalsize\bfseries\color{UdeCAzul}}{\thesubsubsection}{0.6em}{}

% =========================
%        ENCABEZADO
% =========================
\pagestyle{fancy}
\fancyhf{}
\lhead{Electrodinámica Clásica}
\rhead{J.R.}
\cfoot{\thepage\ de \pageref{LastPage}}

\fancypagestyle{plain}{
    \fancyhf{}
    \lhead{Electrodinámica Clásica}
    \rhead{J.R.}
    \cfoot{\thepage\ de \pageref{LastPage}}
}

% =========================
%      ENTORNOS TEÓRICOS
% =========================
\theoremstyle{definition}
\newtheorem{definicion}{Definición}[section]
\newtheorem{ejemplo}[definicion]{Ejemplo}
\newtheorem{observacion}[definicion]{Observación}
\newtheorem{notacion}[definicion]{Notación}

\theoremstyle{plain}
\newtheorem{proposicion}[definicion]{Proposición}
\newtheorem{teorema}[definicion]{Teorema}
\newtheorem{lema}[definicion]{Lema}
\newtheorem{corolario}[definicion]{Corolario}

\numberwithin{equation}{section}

\theoremstyle{remark}
\newtheorem*{comentario}{Comentario}

% =========================
%          CAJAS
% =========================
\newtcolorbox{cajaidea}{
    breakable,
    colback=UdeCGrisClaro,
    colframe=UdeCAzul,
    boxrule=0.5pt,
    arc=2mm,
    left=2.2mm,right=2.2mm,top=1.5mm,bottom=1.5mm
}

\newtcolorbox{cuidado}{
    breakable,
    colback=UdeCDorado!8,
    colframe=UdeCDorado,
    boxrule=0.5pt,
    arc=2mm,
    title=\textbf{Cuidado},
    left=2.2mm,right=2.2mm,top=1.5mm,bottom=1.5mm
}

% =========================
%     COMANDOS ÚTILES
% =========================
\newcommand{\R}{\mathbb{R}}
\newcommand{\C}{\mathbb{C}}
\newcommand{\N}{\mathbb{N}}
\newcommand{\ii}{\mathrm{i}}
\newcommand{\ee}{\mathrm{e}}

\newcommand{\E}{\vec{E}}
\newcommand{\B}{\vec{B}}
\newcommand{\D}{\vec{D}}
\newcommand{\HH}{\vec{H}}
\newcommand{\J}{\vec{J}}
\newcommand{\A}{\vec{A}}
\newcommand{\F}{\vec{F}}
\newcommand{\rvec}{\vec{r}}
\newcommand{\uvec}{\vec{u}}
\newcommand{\p}{\partial}
\newcommand{\epsz}{\varepsilon_0}
\newcommand{\muz}{\mu_0}
\newcommand{\lap}{\nabla^2}
\newcommand{\rot}{\nabla\times}
\newcommand{\diver}{\nabla\cdot}

% =========================
%     DATOS DEL DOCUMENTO
% =========================
\newcommand{\TituloApunte}{Electrodinámica 2}
\newcommand{\SubtituloApunte}{Apunte de estudio basado en clases y en las notas de referencia}
\newcommand{\AutorApunte}{José Ignacio Rosas Sepúlveda}
\newcommand{\Institucion}{Universidad de Concepción}
\newcommand{\Carrera}{Licenciatura en Ciencias Físicas}
\newcommand{\SemestreCurso}{Primer semestre de 2026}
\newcommand{\FechaApunte}{13 de abril de 2026}
